\chapter{Megvalósítás}

\textbf{TODO:} megírni ezt a részt.

\section{Rendszerhívások}

Az implementáció közvetlen rendszerhívásokat használ. Ennek elsődleges oka,
hogy a könyvtár olyan Linux-specifikus funkciókat is elérjen, amelyekhez nem
minden esetben érhető el szabványos vagy hordozható libc wrapper. Ilyen
esetekben a rendszerhívás közvetlen meghívása egyszerűbb és kiszámíthatóbb
megoldást biztosít, mivel az implementáció nem függ a C könyvtár által
biztosított absztrakcióktól.

A könyvtárban definiált rendszerhívás-burkoló függvények nevei és paraméterei
többnyire követik a megfelelő Linux rendszerhívások konvencióit. A libc
függvényekkel való névütközés elkerülése érdekében ezek a függvények
\texttt{sys\_} prefixet kapnak. Például a \texttt{write} rendszerhíváshoz
tartozó wrapper neve \texttt{sys\_write}.

A tényleges rendszerhívások végrehajtását a \texttt{raw\_syscall0},
\ldots, \texttt{raw\_syscall6}, valamint az \texttt{exit\_group}
függvények valósítják meg. Ezek a függvények architektúra-specifikus, GCC/Clang
által támogatott kiterjesztett inline assemblyt használnak, és közvetlenül a
Linux rendszerhívási interfészt hívják meg.

Fontos különbség a libc függvényekhez képest, hogy a nyers Linux
rendszerhívások nem az \texttt{errno} változón keresztül jelzik a hibákat.
Sikertelen hívás esetén a visszatérési érték egy negatív hibakód, amely
általában a $[-4095, -1]$ intervallumba esik. Ennek abszolút értéke felel meg
annak a hibakódnak, amelyet a libc wrapper normál esetben az \texttt{errno}
változóba írna. Mivel az implementáció közvetlen rendszerhívásokat használ, a
hibakódok közvetlenül a visszatérési értékből kezelhetők.

\section{Seccomp szűrők létrehozása}

A könyvtár a rendszerhívások korlátozására a Linux \texttt{seccomp}
mechanizmusának \texttt{SECCOMP\_SET\_MODE\_FILTER} módját használja. Ebben az
üzemmódban a folyamat egy klasszikus BPF programot ad át a kernelnek, amely
minden rendszerhívás előtt lefut. A BPF program a rendszerhívás sorszáma, az
architektúra és a rendszerhívás argumentumai alapján dönt arról, hogy az adott
hívás engedélyezett-e.

Az implementációban a seccomp szabályok \texttt{sock\_filter} elemekből álló
tömbként jelennek meg. Ezekből a tömbökből egy \texttt{sock\_fprog} struktúra
készül, amely tartalmazza a program hosszát és a BPF utasításokra mutató
pointert. Ezt a struktúrát adja át a könyvtár a \texttt{seccomp}
rendszerhívásnak.

A BPF program minden esetben egy közös bevezető résszel kezdődik. Ez először
ellenőrzi, hogy a rendszerhívás az elvárt architektúráról, jelen esetben
\texttt{x86\_64}-ről érkezik-e. Amennyiben az architektúra nem egyezik, a
szűrő azonnal a teljes folyamat leállítását kéri a kerneltől. Ezután a program
betölti a rendszerhívás sorszámát, amely alapján a további szabályok
döntenek az engedélyezésről vagy tiltásról.

Az architektúra ellenőrzése különösen fontos, mert a rendszerhívások sorszámai
architektúránként eltérhetnek. Egy másik architektúrán ugyanaz a szám más
rendszerhívást jelenthetne, ezért a szűrő csak akkor értelmezhető biztonságosan,
ha először meggyőződik arról, hogy a \texttt{seccomp\_data.arch} mező értéke
\texttt{AUDIT\_ARCH\_X86\_64}.

A szabályok létrehozását makrók segítik. Az \texttt{ALLOW\_SYSCALL} makró egy
konkrét rendszerhívást engedélyez, míg az
\texttt{ALLOW\_SYSCALL\_RANGE} egy zárt intervallumba eső rendszerhívásokat
enged át. A szabálylista végére mindig egy alapértelmezett tiltó szabály kerül:
ez vagy a teljes folyamatot állítja le
\texttt{SECCOMP\_RET\_KILL\_PROCESS} visszatérési értékkel, vagy csak az
aktuális szálat \texttt{SECCOMP\_RET\_KILL\_THREAD} használatával. Így a
szűrő alapértelmezetten tiltó jellegű: csak azok a rendszerhívások hajthatók
végre, amelyekre a program explicit engedélyező szabályt tartalmaz.

\begin{lstlisting}[language=C++,caption={Seccomp BPF szabályokat létrehozó segédmakrók}]
#define ALLOW_SYSCALL_RANGE(min, max) \
    BPF_JUMP(BPF_JMP + BPF_JGE + BPF_K, min, 0, 2), \
    BPF_JUMP(BPF_JMP + BPF_JGT + BPF_K, max, 1, 0), \
    BPF_STMT(BPF_RET + BPF_K, SECCOMP_RET_ALLOW)

#define ALLOW_SYSCALL(nr) \
    BPF_JUMP(BPF_JMP + BPF_JEQ + BPF_K, nr, 0, 1), \
    BPF_STMT(BPF_RET + BPF_K, SECCOMP_RET_ALLOW)

#define END_KILL_PROCESS() \
    BPF_STMT(BPF_RET + BPF_K, SECCOMP_RET_KILL_PROCESS)
\end{lstlisting}

A könyvtár
több kisebb szabálykategóriát is definiál. Ilyen például a \texttt{basic},
az \texttt{io} és a \texttt{sandbox} kategória. A \texttt{basic} kategória az
alapvető működéshez szükséges rendszerhívásokat tartalmazza, például a
memóriakezeléshez, jelzéskezeléshez és a folyamat befejezéséhez kapcsolódó
hívásokat. Az \texttt{io} kategória a be- és kimeneti műveletekhez szükséges
rendszerhívásokat engedélyezi, például a \texttt{read}, \texttt{write},
\texttt{pread}, \texttt{pwrite}, \texttt{readv}, \texttt{writev} és
\texttt{sendfile} hívásokat. A \texttt{sandbox} kategória azokat a
rendszerhívásokat tartalmazza, amelyek további sandbox mechanizmusok
bekapcsolásához szükségesek, például a \texttt{seccomp} és a Landlockhoz
kapcsolódó rendszerhívások.

\subsection{A seccomp builder működése}

A szűrők összeállítását a \texttt{SeccompBuilder} osztály végzi. A builder
belsőleg egy \texttt{std::vector<sock\_filter>} tárolót használ, amelybe
egymás után kerülnek be a BPF utasítások. Az inicializálás során először a
közös bevezető rész kerül a szűrő elejére, majd az alapértelmezett
\texttt{basic} szabályok. Az \texttt{init\_no\_defaults} változat ezzel
szemben csak a bevezető részt adja hozzá, és nem engedélyez automatikusan
semmilyen kategóriát.

A felhasználó az \texttt{allow} függvénnyel adhat hozzá további
szabálykategóriákat. A függvény név alapján megkeresi a megfelelő
\texttt{sock\_fprog} struktúrát, majd annak utasításait hozzáfűzi a builder
belső szűrőlistájához. Ha a megadott kategória nem létezik, a függvény
\texttt{ENOENT} hibával tér vissza. Ha a buildert már lezárták, akkor újabb
szabály hozzáadása nem engedélyezett, ilyenkor \texttt{EALREADY} hibát ad
vissza.

A szűrő véglegesítését a \texttt{build\_static} függvény végzi. Ekkor kerül
a szabálylista végére az alapértelmezett tiltó szabály, amely vagy a teljes
folyamatot, vagy csak az aktuális szálat állítja le. Ezután a builder lezárt
állapotba kerül, így a szűrő tartalma többé nem módosítható.

A véglegesített BPF program számára az implementáció külön memóriaterületet
foglal a \texttt{mmap} rendszerhívással. A builder először írható és olvasható
memóriát kér, majd ebbe átmásolja az összeállított \texttt{sock\_filter}
elemeket. A másolás után a memóriaterületet \texttt{mprotect} segítségével
csak olvashatóvá teszi. Ennek célja, hogy a már elkészült szűrőprogram ne
módosulhasson véletlenül a további futás során.

A létrehozott szűrőre a \texttt{SeccompRuleView} objektum mutat. Ez egy
nézet jellegű objektum, amely a \texttt{sock\_fprog} struktúrát tartalmazza.
A \texttt{build} függvény erre építve hoz létre birtokló
\texttt{SeccompRule} objektumot.

A kész szűrő aktiválása a \texttt{seccomp} rendszerhívással történik. Az
implementáció a \texttt{SECCOMP\_SET\_MODE\_FILTER} műveletet használja, és
a korábban létrehozott \texttt{sock\_fprog} struktúrát adja át a kernelnek.
Siker esetén a hívás nulla értékkel tér vissza. Hiba esetén a nyers
rendszerhívás negatív hibakódot ad vissza, amelyet a könyvtár
\texttt{std::expected<void, int>} formában továbbít a hívó felé.
